\begin{figure}[t]
\centering
\setlength{\tabcolsep}{1.6pt}
\renewcommand{\arraystretch}{1.05}
\begin{tabular}{@{}cccccc@{}}
\scriptsize Current frame
& \scriptsize Recorded future
& \scriptsize Recorded change
& \scriptsize ImageWAM
& \scriptsize \modelname{}, offset $[0]$
& \scriptsize \modelname{}, offset $[1]$ \\
\includegraphics[width=0.156\linewidth]{Figures/readout_diff/agent_current.png}
& \includegraphics[width=0.156\linewidth]{Figures/readout_diff/agent_future.png}
& \includegraphics[width=0.156\linewidth]{Figures/readout_diff/agent_true_change_signed.png}
& \includegraphics[width=0.156\linewidth]{Figures/readout_diff/agent_native_signed.png}
& \includegraphics[width=0.156\linewidth]{Figures/readout_diff/agent_jam_off0_signed.png}
& \includegraphics[width=0.156\linewidth]{Figures/readout_diff/agent_jam_signed.png} \\
\multicolumn{6}{c}{\scriptsize Darker \enspace \includegraphics[width=0.22\linewidth,height=0.55em]{Figures/readout_diff/signed_colorbar.png} \enspace Brighter}
\end{tabular}
\vspace{-3mm}
\caption{Decoded future change for \emph{put both moka pots on the stove}, for ImageWAM and the two \modelname{} (Ovis-U1) policies of Table~\ref{tab:objectives} with readout offsets $[0]$ and $[1]$. Two MLP readouts reconstruct the current frame and the frame $16$ steps ahead from the same frozen features, and their difference is compared with the recorded change. Red marks pixels that get brighter and blue marks pixels that get darker, on a common scale. In this example, the features of \modelname{} trained with offset $[1]$ make the future-frame change more readily decodable than the features of ImageWAM or \modelname{} trained with offset $[0]$.}
\label{fig:readoutdiff}
\vspace{-5mm}
\end{figure}
